\begin{bt}%[9D6H18-2]
% nguon: Đề tham khảo Tuyển sinh 10 TP.HCM 2025-2026 – Quận 3 – Đề 1 – Bài 1
Cho parabol $(P): y = \dfrac{1}{2}x^2$.
\begin{enumerate}[a)]
	\item Vẽ đồ thị $(P)$ trên mặt phẳng toạ độ $Oxy$.
	\item Tìm điểm $K$ khác gốc toạ độ thuộc $(P)$ sao cho tung độ gấp đôi hoành độ.
\end{enumerate}
\loigiai{
\begin{enumerate}[a)]
	\item Bảng giá trị tương ứng giữa $x$ và $y$ của hàm số $y = \dfrac{1}{2}x^2$:
	\begin{center}
	\begin{tabular}{|c|c|c|c|c|c|}
	\hline
	$x$ & $-4$ & $-2$ & $0$ & $2$ & $4$ \\
	\hline
	$y = \dfrac{1}{2}x^2$ & $8$ & $2$ & $0$ & $2$ & $8$ \\
	\hline
	\end{tabular}
	\end{center}
	Đồ thị $(P)$ là một parabol đi qua gốc toạ độ $O(0; 0)$, nhận trục tung $Oy$ làm trục đối xứng và nằm phía trên trục hoành.
	\item Điểm $K(x; y)$ thuộc $(P)$ có tung độ gấp đôi hoành độ thoả mãn:
	$y = 2x \Leftrightarrow \dfrac{1}{2}x^2 = 2x \Leftrightarrow x^2 - 4x = 0 \Leftrightarrow x(x - 4) = 0 \Leftrightarrow \left[\begin{aligned} x &= 0 \Rightarrow y = 0 \\ x &= 4 \Rightarrow y = 8. \end{aligned}\right.$\\
	Vì điểm $K$ khác gốc toạ độ nên $x \neq 0$. Do đó $K(4; 8)$.
\end{enumerate}
}
\end{bt}

\begin{bt}%[9D6H20-2]
% nguon: Đề tham khảo Tuyển sinh 10 TP.HCM 2025-2026 – Quận 3 – Đề 1 – Bài 2
Cho phương trình $2x^2 + 4x - 1 = 0$. Không giải phương trình:
\begin{enumerate}[a)]
	\item Chứng tỏ phương trình luôn có hai nghiệm phân biệt.
	\item Tính giá trị của biểu thức $A = \dfrac{1}{2}x_1(2x_1 + 3x_2) + \dfrac{1}{4}x_2(4x_2 + 6x_1)$.
\end{enumerate}
\loigiai{
\begin{enumerate}[a)]
	\item Xét phương trình $2x^2 + 4x - 1 = 0$, ta có hệ số $a = 2$ và $c = -1$.\\
	Vì tích $a\cdot c = 2\cdot(-1) = -2 < 0$ nên phương trình luôn có hai nghiệm phân biệt trái dấu $x_1, x_2$.
	\item Theo định lí Viète, ta có:
	$\begin{cases} x_1 + x_2 = -\dfrac{b}{a} = -\dfrac{4}{2} = -2 \\ x_1 x_2 = \dfrac{c}{a} = -\dfrac{1}{2}. \end{cases}$\\
	Biến đổi biểu thức $A$:
	$A = x_1^2 + \dfrac{3}{2}x_1 x_2 + x_2^2 + \dfrac{3}{2}x_1 x_2 = (x_1^2 + x_2^2) + 3x_1 x_2 = (x_1 + x_2)^2 + x_1 x_2$.\\
	Thay $x_1 + x_2 = -2$ và $x_1 x_2 = -\dfrac{1}{2}$ vào biểu thức $A$, ta được:
	$A = (-2)^2 + \left(-\dfrac{1}{2}\right) = 4 - \dfrac{1}{2} = \dfrac{7}{2} = 3{,}5$.\\
	Vậy $A = \dfrac{7}{2} = 3{,}5$.
\end{enumerate}
}
\end{bt}

\begin{bt}%[9D8H26-3]
% nguon: Đề tham khảo Tuyển sinh 10 TP.HCM 2025-2026 – Quận 3 – Đề 1 – Bài 3
Bạn An đến một hội chợ tham gia trò chơi ném bi vào một bảng điểm có $25$ ô gồm: $5$ ô điểm $5$, $12$ ô điểm $3$, $4$ ô điểm $-2$ và $4$ ô điểm $-1$. Quy định tính điểm và đổi thưởng như sau:
\begin{itemize}
	\item Ném ra ngoài bảng bị trừ $5$ điểm.
	\item Sau $10$ lần ném: Đạt từ $30$ đến $49$ điểm nhận phần quà trị giá $300\,000\text{ đồng}$; đạt từ $15$ đến $29$ điểm nhận phần quà $50\,000\text{ đồng}$.
\end{itemize}
Trong $9$ lần ném đầu tiên, bạn An ném được: $5$ lần vào ô $5$ điểm, $1$ lần ra ngoài bảng, $2$ lần vào ô $3$ điểm và $1$ lần vào ô $-1$ điểm.
\begin{enumerate}[a)]
	\item Tính số điểm bạn An nhận được sau $9$ lần ném.
	\item Hỏi An có cơ hội nhận phần quà trị giá $300\,000\text{ đồng}$ không? Nếu có thì ở lần ném thứ $10$, An phải ném trúng ô mấy điểm? Tính xác suất để An nhận được phần quà này (giả thiết ném trúng các ô trong bảng là đồng khả năng và không ném ra ngoài).
\end{enumerate}
\loigiai{
\begin{enumerate}[a)]
	\item Tổng số điểm của An sau $9$ lần ném là:
	$5 \cdot 5 + 1 \cdot (-5) + 2 \cdot 3 + 1 \cdot (-1) = 25 - 5 + 6 - 1 = 25\text{ (điểm)}$.
	\item Để nhận được phần quà trị giá $300\,000\text{ đồng}$, tổng số điểm sau $10$ lần ném phải đạt từ $30$ đến $49$ điểm.\\
	Gọi số điểm ở lần ném thứ $10$ là $x$, ta có:
	$30 \le 25 + x \le 49 \Leftrightarrow 5 \le x \le 24$.\\
	Trong các ô điểm của bảng ($-2; -1; 3; 5$), chỉ có điểm $5$ thoả mãn điều kiện $5 \le x \le 24$.\\
	Như vậy, An hoàn toàn có cơ hội nhận phần quà $300\,000\text{ đồng}$ nếu ở lần ném thứ $10$ bạn ném trúng ô $5$ điểm.\\
	Vì bảng có tổng cộng $25$ ô và có đúng $5$ ô mang điểm $5$, nên xác suất để An ném trúng ô $5$ điểm là:
	$P = \dfrac{5}{25} = \dfrac{1}{5} = 0{,}2$.
\end{enumerate}
}
\end{bt}

\begin{bt}%[9D1H2-1]
% nguon: Đề tham khảo Tuyển sinh 10 TP.HCM 2025-2026 – Quận 3 – Đề 1 – Bài 4
Một tổ may gồm cả công nhân nam và nữ được giao may $360$ chiếc áo cổ vũ. Để hoàn thành nhiệm vụ, mỗi công nhân nam may $6$ chiếc áo, mỗi công nhân nữ may $8$ chiếc áo. Gọi $x$ và $y$ lần lượt là số công nhân nam và số công nhân nữ ($x, y \in \mathbb{N}^*$).
\begin{enumerate}[a)]
	\item Lập công thức biểu diễn $y$ theo $x$.
	\item Tổ may có bao nhiêu công nhân? Biết rằng số công nhân nữ gấp $1{,}5$ lần số công nhân nam.
\end{enumerate}
\loigiai{
\begin{enumerate}[a)]
	\item Vì $x$ công nhân nam may được $6x$ chiếc áo và $y$ công nhân nữ may được $8y$ chiếc áo, tổng số áo may được là $360$ chiếc nên ta có phương trình:
	$6x + 8y = 360 \Leftrightarrow 8y = 360 - 6x \Leftrightarrow y = \dfrac{360 - 6x}{8} = 45 - \dfrac{3}{4}x$.
	\item Theo giả thiết, số công nhân nữ gấp $1{,}5$ lần số công nhân nam, tức là:
	$y = 1{,}5x = \dfrac{3}{2}x$.\\
	Thay vào phương trình $6x + 8y = 360$, ta được:
	$6x + 8\left(\dfrac{3}{2}x\right) = 360 \Leftrightarrow 6x + 12x = 360 \Leftrightarrow 18x = 360 \Leftrightarrow x = 20\text{ (công nhân)}$.\\
	Suy ra số công nhân nữ là: $y = 1{,}5 \cdot 20 = 30\text{ (công nhân)}$.\\
	Vậy tổ may có tất cả $20 + 30 = 50\text{ công nhân}$.
\end{enumerate}
}
\end{bt}

\begin{bt}%[9H0V31-3]
% nguon: Đề tham khảo Tuyển sinh 10 TP.HCM 2025-2026 – Quận 3 – Đề 1 – Bài 5
Một cửa hàng làm nến thơm dạng hình trụ với lọ thuỷ tinh đựng có chiều cao bên ngoài là $11\text{ cm}$, bán kính đáy bên ngoài là $5\text{ cm}$, độ dày thành lọ là $6\text{ mm}$ ($0{,}6\text{ cm}$) và độ dày đáy lọ là $5\text{ mm}$ ($0{,}5\text{ cm}$).
\begin{enumerate}[a)]
	\item Hỏi lượng sáp thơm cần đổ vào một lọ là bao nhiêu mililit biết cửa hàng cần chừa lại $8\%$ không gian trống bên trong lọ? (Làm tròn kết quả đến hàng phần trăm của $\text{ml}$).
	\item Cửa hàng cần nhập các miếng sáp nguyên chất dạng hình hộp chữ nhật kích thước $30\text{ cm} \times 30\text{ cm} \times 1\text{ cm}$ để sản xuất $200$ lọ nến thơm như trên. Hỏi cửa hàng cần đặt ít nhất bao nhiêu miếng sáp nguyên chất?
\end{enumerate}
\loigiai{
\begin{enumerate}[a)]
	\item Kích thước lòng trong của chiếc lọ hình trụ là:
	\begin{itemize}
		\item Bán kính đáy lòng trong: $r = 5 - 0{,}6 = 4{,}4\text{ (cm)}$.
		\item Chiều cao lòng trong: $h = 11 - 0{,}5 = 10{,}5\text{ (cm)}$.
	\end{itemize}
	Thể tích lòng trong của chiếc lọ là:
	$V_{\text{trong}} = \pi r^2 h = \pi \cdot (4{,}4)^2 \cdot 10{,}5 = 203{,}28\pi\text{ (cm}^3\text{)}$.\\
	Vì cần chừa $8\%$ không gian trống nên lượng sáp thơm đổ vào mỗi lọ chiếm $92\%$ thể tích lòng trong:
	$V_{\text{sáp}} = 92\% \cdot V_{\text{trong}} = 0{,}92 \cdot 203{,}28\pi = 187{,}0176\pi \approx 587{,}53\text{ (cm}^3) = 587{,}53\text{ (ml)}$.
	\item Thể tích một miếng sáp nguyên chất hình hộp chữ nhật là:
	$V_{\text{miếng}} = 30 \cdot 30 \cdot 1 = 900\text{ (cm}^3\text{)}$.\\
	Tổng thể tích sáp thơm cần dùng để làm $200$ lọ nến là:
	$V_{\text{tổng}} = 200 \cdot V_{\text{sáp}} \approx 200 \cdot 587{,}533 = 117\,506{,}6\text{ (cm}^3\text{)}$.\\
	Số miếng sáp nguyên chất cần đặt ít nhất là:
	$n = \dfrac{V_{\text{tổng}}}{V_{\text{miếng}}} = \dfrac{117\,506{,}6}{900} \approx 130{,}56$.\\
	Vì số miếng sáp là số nguyên nên cửa hàng cần đặt ít nhất $131$ miếng sáp nguyên chất.
\end{enumerate}
}
\end{bt}

\begin{bt}%[9D1V3-4]
% nguon: Đề tham khảo Tuyển sinh 10 TP.HCM 2025-2026 – Quận 3 – Đề 1 – Bài 6
Ba chiếc bình chứa nước có tổng dung tích là $105\text{ lít}$. Nếu đổ đầy nước vào bình thứ nhất rồi lấy toàn bộ lượng nước đó đem rót vào hai bình còn lại thì:
\begin{itemize}
	\item Hoặc làm đầy bình thứ hai và rót được thêm $\dfrac{4}{5}$ dung tích bình thứ ba.
	\item Hoặc làm đầy bình thứ ba và rót được thêm $\dfrac{5}{6}$ dung tích bình thứ hai.
\end{itemize}
Hãy xác định thể tích của mỗi bình (giả sử nước không bị hao hụt trong quá trình rót).
\loigiai{
Gọi dung tích của bình thứ nhất, bình thứ hai và bình thứ ba lần lượt là $x, y, z$ (lít; điều kiện: $x, y, z > 0$).\\
Theo đề bài, tổng dung tích ba bình là $105\text{ lít}$, nên ta có:
$x + y + z = 105$.\\
Từ hai trường hợp rót nước của bình thứ nhất, ta có hệ phương trình:
$\begin{cases} x = y + \dfrac{4}{5}z \\ x = z + \dfrac{5}{6}y. \end{cases}$\\
Suy ra:
$y + \dfrac{4}{5}z = z + \dfrac{5}{6}y \Leftrightarrow y - \dfrac{5}{6}y = z - \dfrac{4}{5}z \Leftrightarrow \dfrac{1}{6}y = \dfrac{1}{5}z \Leftrightarrow z = \dfrac{5}{6}y$.\\
Thay $z = \dfrac{5}{6}y$ vào biểu thức của $x$, ta được:
$x = y + \dfrac{4}{5}\left(\dfrac{5}{6}y\right) = y + \dfrac{2}{3}y = \dfrac{5}{3}y$.\\
Thay $x = \dfrac{5}{3}y$ và $z = \dfrac{5}{6}y$ vào phương trình tổng dung tích, ta có:
$\dfrac{5}{3}y + y + \dfrac{5}{6}y = 105 \Leftrightarrow \left(\dfrac{10 + 6 + 5}{6}\right)y = 105 \Leftrightarrow \dfrac{21}{6}y = 105 \Leftrightarrow \dfrac{7}{2}y = 105 \Leftrightarrow y = 30\text{ (lít)}$.\\
Từ đó suy ra:
\begin{itemize}
	\item $x = \dfrac{5}{3} \cdot 30 = 50\text{ (lít)}$.
	\item $z = \dfrac{5}{6} \cdot 30 = 25\text{ (lít)}$.
\end{itemize}
Vậy dung tích của bình thứ nhất là $50\text{ lít}$, bình thứ hai là $30\text{ lít}$ và bình thứ ba là $25\text{ lít}$.
}
\end{bt}

\begin{bt}%[9H9C29-3]
% nguon: Đề tham khảo Tuyển sinh 10 TP.HCM 2025-2026 – Quận 3 – Đề 1 – Bài 7
Cho đường tròn $(O; R)$ và điểm $A$ nằm ngoài $(O)$ sao cho $OA = 2R$. Từ $A$ kẻ hai tiếp tuyến $AB, AC$ đến đường tròn $(O)$ ($B, C$ là hai tiếp điểm). Gọi $H$ là giao điểm của $OA$ và $BC$.
\begin{enumerate}[a)]
	\item Chứng minh tứ giác $ABOC$ nội tiếp đường tròn.
	\item Chứng minh $OA \perp BC$ tại $H$ và $BC^2 = 4 OH \cdot HA$.
	\item Qua $O$ vẽ đường thẳng vuông góc với $OA$ cắt đường thẳng $AB$ tại $D$. Gọi $Q$ là giao điểm của $DH$ và $OB$. Tính diện tích tam giác $QOA$ theo $R$.
\end{enumerate}
\loigiai{
\begin{enumerate}[a)]
	\item Vì $AB, AC$ là các tiếp tuyến của $(O)$ nên $\widehat{ABO} = 90^\circ$ và $\widehat{ACO} = 90^\circ$.\\
	Xét tứ giác $ABOC$, ta có $\widehat{ABO} + \widehat{ACO} = 90^\circ + 90^\circ = 180^\circ$.\\
	Vậy tứ giác $ABOC$ nội tiếp đường tròn đường kính $OA$.
	\item Vì $AB = AC$ (tính chất hai tiếp tuyến cắt nhau) và $OB = OC = R$ nên $OA$ là đường trung trực của đoạn thẳng $BC$, suy ra $OA \perp BC$ tại $H$ và $H$ là trung điểm của $BC \Rightarrow BC = 2BH$.\\
	Xét tam giác $ABO$ vuông tại $B$ có đường cao $BH$. Áp dụng hệ thức lượng trong tam giác vuông:
	$BH^2 = OH \cdot HA$.\\
	Do đó: $BC^2 = (2BH)^2 = 4BH^2 = 4 OH \cdot HA$.
	\item Trong tam giác vuông $ABO$ tại $B$ có $OA = 2R$ và $OB = R$, ta có:
	$\cos\widehat{AOB} = \dfrac{OB}{OA} = \dfrac{R}{2R} = \dfrac{1}{2} \Rightarrow \widehat{AOB} = 60^\circ \Rightarrow \widehat{OAB} = 30^\circ$.\\
	Xét tam giác $OAD$ vuông tại $O$ (do $OD \perp OA$), ta có:
	$OD = OA \cdot \tan\widehat{OAD} = 2R \cdot \tan 30^\circ = 2R \cdot \dfrac{\sqrt{3}}{3} = \dfrac{2R\sqrt{3}}{3}$.\\
	Ta có $OH = OB \cdot \cos 60^\circ = \dfrac{R}{2}$ và $AH = OA - OH = \dfrac{3R}{2}$.\\
	Đặt hệ toạ độ phẳng $Oxy$ với gốc tại $O$, trục $Ox$ hướng theo tia $OA$, trục $Oy$ hướng theo tia $OD$ ($Oy \perp Ox$):
	Khi đó toạ độ các điểm là: $O(0; 0)$, $A(2R; 0)$, $H\left(\dfrac{R}{2}; 0\right)$ và $D\left(0; \dfrac{2R\sqrt{3}}{3}\right)$.\\
	Phương trình đường thẳng $DH$ đi qua hai điểm $D$ và $H$ là:
	$\dfrac{x}{\frac{R}{2}} + \dfrac{y}{\frac{2R\sqrt{3}}{3}} = 1 \Leftrightarrow \dfrac{2x}{R} + \dfrac{3y}{2R\sqrt{3}} = 1 \Leftrightarrow 4x + \sqrt{3}y = 2R$.\\
	Đường thẳng $OB$ đi qua gốc toạ độ $O$ tạo với trục $Ox$ một góc $60^\circ$ nên có phương trình:
	$y = \tan 60^\circ \cdot x \Leftrightarrow y = x\sqrt{3}$.\\
	Giao điểm $Q$ của $DH$ và $OB$ có toạ độ thoả mãn hệ phương trình:
	$\begin{cases} 4x + \sqrt{3}y = 2R \\ y = x\sqrt{3} \end{cases} \Leftrightarrow \begin{cases} 4x + 3x = 2R \\ y = x\sqrt{3} \end{cases} \Leftrightarrow \begin{cases} x_Q = \dfrac{2}{7}R \\ y_Q = \dfrac{2\sqrt{3}}{7}R. \end{cases}$\\
	Tam giác $QOA$ có cạnh đáy $OA$ nằm trên trục $Ox$ với độ dài $OA = 2R$, chiều cao tương ứng là tung độ của $Q$: $h = y_Q = \dfrac{2\sqrt{3}}{7}R$.\\
	Diện tích tam giác $QOA$ là:
	$S_{\Delta QOA} = \dfrac{1}{2} OA \cdot y_Q = \dfrac{1}{2} \cdot 2R \cdot \dfrac{2\sqrt{3}}{7}R = \dfrac{2\sqrt{3}}{7}R^2$.
\end{enumerate}
}
\end{bt}

\begin{bt}%[9D6H18-2]
% nguon: Đề tham khảo Tuyển sinh 10 TP.HCM 2025-2026 – Quận 3 – Đề 2 – Bài 1
Cho parabol $(P): y = -\dfrac{x^2}{4}$.
\begin{enumerate}[a)]
	\item Vẽ đồ thị $(P)$ trên hệ trục toạ độ $Oxy$.
	\item Tìm toạ độ những điểm $M$ thuộc $(P)$ có tung độ gấp $3$ lần hoành độ.
\end{enumerate}
\loigiai{
\begin{enumerate}[a)]
	\item Bảng giá trị tương ứng giữa $x$ và $y$ của hàm số $y = -\dfrac{x^2}{4}$:
	\begin{center}
	\begin{tabular}{|c|c|c|c|c|c|}
	\hline
	$x$ & $-4$ & $-2$ & $0$ & $2$ & $4$ \\
	\hline
	$y = -\dfrac{x^2}{4}$ & $-4$ & $-1$ & $0$ & $-1$ & $-4$ \\
	\hline
	\end{tabular}
	\end{center}
	Đồ thị $(P)$ là một parabol có đỉnh tại gốc toạ độ $O(0; 0)$, nhận trục tung $Oy$ làm trục đối xứng và nằm phía dưới trục hoành.
	\item Điểm $M(x; y)$ thuộc $(P)$ có tung độ gấp $3$ lần hoành độ thoả mãn:
	$y = 3x \Leftrightarrow -\dfrac{x^2}{4} = 3x \Leftrightarrow x^2 + 12x = 0 \Leftrightarrow x(x + 12) = 0 \Leftrightarrow \left[\begin{aligned} x &= 0 \Rightarrow y = 0 \\ x &= -12 \Rightarrow y = -36. \end{aligned}\right.$\\
	Vậy có hai điểm thoả mãn yêu cầu đề bài là $M_1(0; 0)$ và $M_2(-12; -36)$.
\end{enumerate}
}
\end{bt}

\begin{bt}%[9D6H20-2]
% nguon: Đề tham khảo Tuyển sinh 10 TP.HCM 2025-2026 – Quận 3 – Đề 2 – Bài 2
Cho phương trình $3x^2 - 2x - 2 = 0$.
\begin{enumerate}[a)]
	\item Chứng tỏ phương trình luôn có hai nghiệm phân biệt $x_1, x_2$.
	\item Không giải phương trình, hãy tính giá trị của biểu thức $A = x_1^3 + x_2^3 - 3x_1^2 x_2 - 3x_1 x_2^2$.
\end{enumerate}
\loigiai{
\begin{enumerate}[a)]
	\item Xét phương trình $3x^2 - 2x - 2 = 0$, ta có:
	$a\cdot c = 3\cdot(-2) = -6 < 0$.\\
	Vì tích $ac < 0$ nên phương trình luôn có hai nghiệm phân biệt trái dấu $x_1, x_2$.
	\item Theo định lí Viète, ta có:
	$\begin{cases} x_1 + x_2 = -\dfrac{b}{a} = \dfrac{2}{3} \\ x_1 x_2 = \dfrac{c}{a} = -\dfrac{2}{3}. \end{cases}$\\
	Biến đổi biểu thức $A$:
	$A = (x_1^3 + x_2^3) - 3x_1 x_2 (x_1 + x_2) = (x_1 + x_2)(x_1^2 - x_1 x_2 + x_2^2) - 3x_1 x_2(x_1 + x_2) = (x_1 + x_2)\left[(x_1 + x_2)^2 - 3x_1 x_2\right] - 3x_1 x_2(x_1 + x_2) = (x_1 + x_2)\left[(x_1 + x_2)^2 - 6x_1 x_2\right]$.\\
	Thay $x_1 + x_2 = \dfrac{2}{3}$ và $x_1 x_2 = -\dfrac{2}{3}$ vào biểu thức $A$, ta được:
	$A = \dfrac{2}{3}\left[\left(\dfrac{2}{3}\right)^2 - 6\cdot\left(-\dfrac{2}{3}\right)\right] = \dfrac{2}{3}\left(\dfrac{4}{9} + 4\right) = \dfrac{2}{3} \cdot \dfrac{40}{9} = \dfrac{80}{27}$.\\
	Vậy $A = \dfrac{80}{27}$.
\end{enumerate}
}
\end{bt}

\begin{bt}%[9D8H26-2]
% nguon: Đề tham khảo Tuyển sinh 10 TP.HCM 2025-2026 – Quận 3 – Đề 2 – Bài 3
Nhân dịp tổng kết năm học, lớp 9A có $2$ bạn nam hát hay là Tuấn, Hùng và $2$ bạn nữ hát hay là Lan, Hồng. Giáo viên chủ nhiệm chọn ngẫu nhiên $2$ bạn để hát song ca.
\begin{enumerate}[a)]
	\item Hãy liệt kê tất cả các cách chọn ngẫu nhiên $2$ bạn để hát song ca.
	\item Tính xác suất của mỗi biến cố sau:
	\begin{itemize}
		\item Biến cố $A$: ``Trong $2$ bạn được chọn có đúng $1$ bạn nam và $1$ bạn nữ''.
		\item Biến cố $B$: ``Trong $2$ bạn được chọn có bạn Lan''.
	\end{itemize}
\end{enumerate}
\loigiai{
\begin{enumerate}[a)]
	\item Tập hợp tất cả các cặp gồm $2$ bạn được chọn từ $4$ bạn (không phân biệt thứ tự) là không gian mẫu:
	$\Omega = \{(\text{Tuấn}, \text{Hùng}); (\text{Tuấn}, \text{Lan}); (\text{Tuấn}, \text{Hồng}); (\text{Hùng}, \text{Lan}); (\text{Hùng}, \text{Hồng}); (\text{Lan}, \text{Hồng})\}$.\\
	Số phần tử của không gian mẫu là $n(\Omega) = 6$.
	\item Xác suất của từng biến cố:
	\begin{itemize}
		\item Các kết quả thuận lợi cho biến cố $A$ (gồm 1 nam và 1 nữ) là:
		$(\text{Tuấn}, \text{Lan})$, $(\text{Tuấn}, \text{Hồng})$, $(\text{Hùng}, \text{Lan})$ và $(\text{Hùng}, \text{Hồng}) \Rightarrow n(A) = 4$.\\
		Do đó xác suất của biến cố $A$ là:
		$P(A) = \dfrac{n(A)}{n(\Omega)} = \dfrac{4}{6} = \dfrac{2}{3}$.
		\item Các kết quả thuận lợi cho biến cố $B$ (có bạn Lan) là:
		$(\text{Tuấn}, \text{Lan})$, $(\text{Hùng}, \text{Lan})$ và $(\text{Lan}, \text{Hồng}) \Rightarrow n(B) = 3$.\\
		Do đó xác suất của biến cố $B$ là:
		$P(B) = \dfrac{n(B)}{n(\Omega)} = \dfrac{3}{6} = \dfrac{1}{2}$.
	\end{itemize}
\end{enumerate}
}
\end{bt}

\begin{bt}%[9H0V31-1]
% nguon: Đề tham khảo Tuyển sinh 10 TP.HCM 2025-2026 – Quận 3 – Đề 2 – Bài 4
Một hình chóp tứ giác đều có chiều dài cạnh đáy là $x\text{ (cm)}$ và chiều dài đường cao của mặt bên (trung đoạn) ngắn hơn cạnh đáy $2\text{ cm}$.
\begin{enumerate}[a)]
	\item Viết biểu thức $S_{xq}$ biểu diễn diện tích xung quanh của hình chóp theo $x$.
	\item Biết diện tích toàn phần của hình chóp tứ giác đều này là $207\text{ cm}^2$, hãy tính độ dài cạnh đáy $x$.
\end{enumerate}
\loigiai{
\begin{enumerate}[a)]
	\item Chiều dài đường cao của mặt bên (trung đoạn của hình chóp) là: $d = x - 2\text{ (cm)}$ (với $x > 2$).\\
	Hình chóp tứ giác đều có $4$ mặt bên là các tam giác cân có đáy $x$ và đường cao $d = x - 2$.\\
	Biểu thức biểu diễn diện tích xung quanh của hình chóp là:
	$S_{xq} = 4 \cdot \left(\dfrac{1}{2} x \cdot (x - 2)\right) = 2x(x - 2) = 2x^2 - 4x\text{ (cm}^2\text{)}$.
	\item Diện tích đáy của hình chóp tứ giác đều là: $S_{\text{đáy}} = x^2\text{ (cm}^2\text{)}$.\\
	Diện tích toàn phần của hình chóp là:
	$S_{tp} = S_{xq} + S_{\text{đáy}} = (2x^2 - 4x) + x^2 = 3x^2 - 4x\text{ (cm}^2\text{)}$.\\
	Theo đề bài, $S_{tp} = 207\text{ cm}^2$, ta có phương trình:
	$3x^2 - 4x = 207 \Leftrightarrow 3x^2 - 4x - 207 = 0$.\\
	Ta có $\Delta' = (-2)^2 - 3\cdot(-207) = 4 + 621 = 625 = 25^2 > 0$.\\
	Vì $x > 2$ nên phương trình có nghiệm hợp lệ duy nhất:
	$x = \dfrac{2 + 25}{3} = \dfrac{27}{3} = 9\text{ (cm)}$.\\
	Vậy độ dài cạnh đáy của hình chóp tứ giác đều là $9\text{ cm}$.
\end{enumerate}
}
\end{bt}

\begin{bt}%[9H0V31-2]
% nguon: Đề tham khảo Tuyển sinh 10 TP.HCM 2025-2026 – Quận 3 – Đề 2 – Bài 5
Để làm một chiếc hộp đựng bỏng ngô hình nón, từ một miếng giấy hình tròn bán kính $R = 25\text{ cm}$, bạn An cắt bỏ phần hình quạt tròn có góc ở tâm là $135^\circ$, sau đó dán phần hình quạt lớn còn lại để tạo thành chiếc hộp hình nón.
\begin{enumerate}[a)]
	\item Tính diện tích giấy được sử dụng để làm chiếc hộp trên (làm tròn đến hàng phần trăm của $\text{cm}^2$).
	\item Biết chiếc hộp có thể chứa nhiều hơn thể tích thực của nó tối đa là $15\%$ khi bỏng ngô vun lên phía trên miệng hộp. Tính thể tích bỏng ngô tối đa mà chiếc hộp có thể chứa được (làm tròn kết quả đến hàng đơn vị của $\text{cm}^3$).
\end{enumerate}
\loigiai{
\begin{enumerate}[a)]
	\item Góc ở tâm của phần hình quạt lớn dùng làm hộp nón là:
	$\alpha = 360^\circ - 135^\circ = 225^\circ$.\\
	Diện tích giấy sử dụng để làm hộp chính là diện tích của hình quạt tròn này:
	$S = \dfrac{\pi R^2 \cdot 225^\circ}{360^\circ} = \dfrac{5}{8}\pi \cdot 25^2 = \dfrac{3125\pi}{8} \approx 1227{,}18\text{ (cm}^2\text{)}$.
	\item Độ dài đường sinh của hình nón bằng bán kính của hình quạt tròn: $l = R = 25\text{ cm}$.\\
	Chu vi đáy của hình nón bằng độ dài cung của hình quạt tròn góc $225^\circ$:
	$C = \dfrac{2\pi R \cdot 225^\circ}{360^\circ} = \dfrac{5}{8} \cdot 2\pi \cdot 25 = \dfrac{125\pi}{4}\text{ (cm)}$.\\
	Bán kính đáy $r$ của hình nón thoả mãn: $2\pi r = \dfrac{125\pi}{4} \Rightarrow r = \dfrac{125}{8} = 15{,}625\text{ (cm)}$.\\
	Chiều cao của hình nón là:
	$h = \sqrt{l^2 - r^2} = \sqrt{25^2 - (15{,}625)^2} = \sqrt{625 - 244{,}140625} = \sqrt{380{,}859375} \approx 19{,}5156\text{ (cm)}$.\\
	Thể tích thực của chiếc hộp hình nón là:
	$V_{\text{thực}} = \dfrac{1}{3}\pi r^2 h = \dfrac{1}{3}\pi \cdot (15{,}625)^2 \cdot 19{,}5156 \approx 4989{,}43\text{ (cm}^3\text{)}$.\\
	Thể tích bỏng ngô tối đa chiếc hộp có thể chứa được (nhiều hơn $15\%$ thể tích thực) là:
	$V_{\max} = 115\% \cdot V_{\text{thực}} = 1{,}15 \cdot 4989{,}43 \approx 5737{,}85\text{ (cm}^3\text{)}$.\\
	Làm tròn đến hàng đơn vị, thể tích bỏng ngô chứa được nhiều nhất là khoảng $5738\text{ cm}^3$.
\end{enumerate}
}
\end{bt}

\begin{bt}%[9D1V3-2]
% nguon: Đề tham khảo Tuyển sinh 10 TP.HCM 2025-2026 – Quận 3 – Đề 2 – Bài 6
Nhân dịp ngày Nhà giáo Việt Nam, nhà trường tổ chức cuộc thi kéo co. Giáo viên chủ nhiệm lớp 9B đã chọn ra $\dfrac{1}{4}$ số học sinh nam và $\dfrac{1}{3}$ số học sinh nữ của lớp tham gia kéo co, biết rằng số học sinh nam và số học sinh nữ tham gia kéo co là bằng nhau. Toàn bộ $20$ học sinh còn lại của lớp tham gia đội cổ động. Hãy tìm số học sinh nam và số học sinh nữ của lớp 9B.
\loigiai{
Gọi số học sinh nam và số học sinh nữ của lớp 9B lần lượt là $x$ và $y$ ($x, y \in \mathbb{N}^*$, $x$ chia hết cho $4$ và $y$ chia hết cho $3$).\\
Số học sinh nam tham gia kéo co là $\dfrac{1}{4}x$ (học sinh).\\
Số học sinh nữ tham gia kéo co là $\dfrac{1}{3}y$ (học sinh).\\
Vì số học sinh nam và số học sinh nữ tham gia kéo co bằng nhau nên ta có phương trình:
$\dfrac{1}{4}x = \dfrac{1}{3}y \Leftrightarrow 3x - 4y = 0$.\\
Số học sinh nam làm cổ động viên là: $x - \dfrac{1}{4}x = \dfrac{3}{4}x$ (học sinh).\\
Số học sinh nữ làm cổ động viên là: $y - \dfrac{1}{3}y = \dfrac{2}{3}y$ (học sinh).\\
Vì đội cổ động có tổng cộng $20$ học sinh nên ta có phương trình:
$\dfrac{3}{4}x + \dfrac{2}{3}y = 20 \Leftrightarrow 9x + 8y = 240$.\\
Ta có hệ phương trình:
$\begin{cases} 3x - 4y = 0 \\ 9x + 8y = 240 \end{cases} \Leftrightarrow \begin{cases} 6x - 8y = 0 \\ 9x + 8y = 240 \end{cases} \Leftrightarrow \begin{cases} 15x = 240 \\ y = \dfrac{3}{4}x \end{cases} \Leftrightarrow \begin{cases} x = 16 \\ y = 12 \end{cases}$ (thoả mãn điều kiện).\\
Vậy lớp 9B có $16$ học sinh nam và $12$ học sinh nữ.
}
\end{bt}

\begin{bt}%[9H9C29-3]
% nguon: Đề tham khảo Tuyển sinh 10 TP.HCM 2025-2026 – Quận 3 – Đề 2 – Bài 7
Cho đường tròn $(O)$ có đường kính $BC$. Trên nửa đường tròn $(O)$, lấy hai điểm $A$ và $D$ theo thứ tự $B, A, D, C$. Tia $BA$ và tia $CD$ cắt nhau tại điểm $S$, đoạn thẳng $AC$ cắt $BD$ tại $H$. Gọi $I$ là trung điểm của đoạn thẳng $SH$.
\begin{enumerate}[a)]
	\item Chứng minh tam giác $ABC$ và tam giác $BDC$ vuông. Chứng minh tứ giác $SAHD$ nội tiếp đường tròn.
	\item Tia $SH$ cắt cạnh $BC$ tại $M$, đường thẳng $DM$ cắt $HC$ tại $K$. Chứng minh $\widehat{IAH} = \widehat{KDC}$.
	\item Trong trường hợp $\widehat{BSC} = 60^\circ$ và $BC = 6\text{ cm}$, tính diện tích hình viên phân giới hạn bởi cung nhỏ $AD$ và dây cung $AD$.
\end{enumerate}
\loigiai{
\begin{enumerate}[a)]
	\item Các góc $\widehat{BAC}$ và $\widehat{BDC}$ là các góc nội tiếp chắn nửa đường tròn $(O)$ đường kính $BC$ nên:
	$\widehat{BAC} = 90^\circ \Rightarrow CA \perp SB$ tại $A$;\\
	$\widehat{BDC} = 90^\circ \Rightarrow BD \perp SC$ tại $D$.\\
	Suy ra các tam giác $\Delta ABC$ và $\Delta BDC$ là các tam giác vuông.\\
	Trong tam giác $SBC$, hai đường cao $CA$ và $BD$ cắt nhau tại $H$, do đó $H$ là trực tâm của $\Delta SBC$.\\
	Xét tứ giác $SAHD$, ta có:
	$\widehat{SAH} = 90^\circ$ (do $CA \perp SB$) và $\widehat{SDH} = 90^\circ$ (do $BD \perp SC$).\\
	Suy ra $\widehat{SAH} + \widehat{SDH} = 90^\circ + 90^\circ = 180^\circ$.\\
	Vậy tứ giác $SAHD$ nội tiếp đường tròn đường kính $SH$ có tâm là trung điểm $I$ của $SH$.
	\item Vì $H$ là trực tâm tam giác $SBC$ nên $SH \perp BC$ tại $M$, suy ra $\widehat{SM C} = 90^\circ$.\\
	Tứ giác $SDHM$ có $\widehat{SDH} = 90^\circ$ và $\widehat{SMH} = 90^\circ \Rightarrow SDHM$ nội tiếp đường tròn đường kính $SH$.\\
	Do đó 5 điểm $S, A, H, D, M$ cùng thuộc đường tròn đường kính $SH$ tâm $I$.\\
	Trong đường tròn $(I)$ đường kính $SH$, tam giác $IAH$ cân tại $I$ (do $IA = IH = \dfrac{SH}{2}$) nên:
	$\widehat{IAH} = \widehat{IHA} = \widehat{SH A}$.\\
	Mặt khác, tứ giác $HDMC$ nội tiếp (do $\widehat{HDC} = 90^\circ$ và $\widehat{HMC} = 90^\circ$), nên:
	$\widehat{KDC} = \widehat{HDC} - \widehat{HDM} = 90^\circ - \widehat{HDM}$.\\
	Trong tứ giác nội tiếp $SDHM$, ta có $\widehat{HDM} = \widehat{HSM}$. Trong tam giác vuông $SH A$ vuông tại $A$, ta có $\widehat{HSM} = 90^\circ - \widehat{SHA}$.\\
	Từ đó suy ra $\widehat{KDC} = \widehat{SHA} = \widehat{IAH}$.
	\item Xét tam giác vuông $SBC$ có $\widehat{BSC} = 60^\circ$:\\
	Trong tam giác vuông $SCA$ tại $A$: $\cos 60^\circ = \dfrac{SA}{SC} = \dfrac{1}{2}$.\\
	Trong tam giác vuông $SBD$ tại $D$: $\cos 60^\circ = \dfrac{SD}{SB} = \dfrac{1}{2}$.\\
	Suy ra $\dfrac{SA}{SC} = \dfrac{SD}{SB} = \dfrac{1}{2}$.\\
	Xét tam giác $SAD$ và tam giác $SCB$, có góc $\widehat{BSC} = 60^\circ$ chung và $\dfrac{SA}{SC} = \dfrac{SD}{SB}$, do đó $\Delta SAD \backsim \Delta SCB$ (c.g.c) theo tỉ số đồng dạng $k = \dfrac{1}{2}$.\\
	Suy ra: $AD = \dfrac{1}{2} BC = \dfrac{1}{2} \cdot 6 = 3\text{ (cm)}$.\\
	Bán kính đường tròn $(O)$ là $R = \dfrac{BC}{2} = \dfrac{6}{2} = 3\text{ cm}$.\\
	Xét tam giác $OAD$ có $OA = OD = R = 3\text{ cm}$ và $AD = 3\text{ cm}$, do đó tam giác $OAD$ là tam giác đều cạnh $3\text{ cm}$.\\
	Suy ra góc ở tâm $\widehat{AOD} = 60^\circ$.\\
	Diện tích hình quạt tròn $OAD$ là:
	$S_{\text{quạt}} = \dfrac{\pi R^2 \cdot 60^\circ}{360^\circ} = \dfrac{\pi \cdot 3^2}{6} = \dfrac{3\pi}{2}\text{ (cm}^2\text{)}$.\\
	Diện tích tam giác đều $OAD$ cạnh $R = 3\text{ cm}$ là:
	$S_{\Delta OAD} = \dfrac{R^2\sqrt{3}}{4} = \dfrac{3^2\sqrt{3}}{4} = \dfrac{9\sqrt{3}}{4}\text{ (cm}^2\text{)}$.\\
	Diện tích hình viên phân giới hạn bởi cung nhỏ $AD$ và dây cung $AD$ là:
	$S_{\text{vp}} = S_{\text{quạt}} - S_{\Delta OAD} = \dfrac{3\pi}{2} - \dfrac{9\sqrt{3}}{4} \approx 4{,}7124 - 3{,}8971 = 0{,}8153\text{ (cm}^2\text{)}$.\\
	Vậy diện tích hình viên phân cần tìm là $\dfrac{3\pi}{2} - \dfrac{9\sqrt{3}}{4} \approx 0{,}82\text{ cm}^2$.
\end{enumerate}
}
\end{bt}
